\section{Intermediate biases}\label{sec:intermediate}

We retain the notation of Section~\ref{sec:continuation}. In particular,
\(21/25\le s<1\), \(r=1-s^2\), \(d=T_\rho f\),
\(m=\E f\ge0\), \(h=H(m)\), \(b=\E H(d)\), and
\(g=F(d)\), where \(F(0)=0\) and \(F'=H^{-3/2}\).
The first-level bound of Lemma~\ref{lem:fourier-bounds} is denoted by
\[
 w=\min\{2p,4\sqrt{3/2}\,p^{3/2}\},\qquad p=(1-m)/2.
\]
All terminating decimals in this section denote exact rational numbers.

\begin{proposition}[Exclusion of an intermediate-bias crossing]
\label{prop:intermediate}
Let \(f\) be a nonconstant Boolean function and suppose
\[
 \frac{r^2}{2}\le m\le\frac{37}{50},\qquad \frac{21}{25}\le s<1.
\]
If \(b=h-1+s\), then \(\mathcal E(d)>r/s\).
Consequently a crossing as in Lemma~\ref{lem:crossing} cannot occur
in this range.
\end{proposition}

\subsection{Parity and an energy reduction}

For a function \(u\) on the cube, write
\[
 u_e(x)=\frac{u(x)+u(-x)}2,\qquad
 u_o(x)=\frac{u(x)-u(-x)}2.
\]
These are the projections onto the even and odd Fourier degrees. Every
nonconstant even degree is at least two, whereas an odd degree other
than one is at least three. As in Lemma~\ref{lem:fourier-bounds},
Parseval and \(W_1(f)\le w\) give
\begin{equation}\label{mid:odd-cap}
 \Var(d_o)
 \le rW_1(f)+r^3\bigl(h^2-W_1(f)\bigr)
 \le r\bigl[(1-r^2)w+r^2h^2\bigr].
\end{equation}
The edge comparison of Lemma~\ref{lem:basic-energy}, followed by the
same degree count for \(g\), gives
\begin{equation}\label{mid:parity-energy}
 \mathcal E(d)\ge\sum_i\E(D_i g)^2
 \ge\Var(g)+\Var(g_e).
\end{equation}

To use the extra even-degree contribution, we compare \(g=F(d)\)
with its affine approximation at \(m\). Set \(v=1-s\) and
\(a=F'(m)=h^{-3/2}\). The tangent deficit of \(H\) at \(m\) is
\begin{equation}\label{mid:V}
 V=h-H(d)-\frac{m(d-m)}h.
\end{equation}
Concavity gives \(V\ge0\), and at a putative crossing,
\begin{equation}\label{mid:EV}
 \E V=h-b=v.
\end{equation}
The variance expansion
\[
 \Var(g)=\Var(g-ad)
       +2a\E\bigl[(d-m)(g-F(m))\bigr]-a^2\Var(d)
\]
is valid without any assumption on \(\E g\). On the nonconstant even
subspace, apply the identity
\[
 \|u-v_0\|_2^2+\|u\|_2^2
 =2\|u-v_0/2\|_2^2+\|v_0\|_2^2/2
\]
with \(u=g_e-\E g\) and \(v_0=a(d_e-m)\). Dropping the odd part
of \(\Var(g-ad)\) then gives
\[
 \Var(g-ad)+\Var(g_e)\ge\frac{a^2}{2}\Var(d_e).
\]
This compensates for half the even part of the variance penalty.
Define the remaining pointwise expression by
\begin{equation}\label{mid:Psi}
 \Psi=2a(d-m)\bigl(g-F(m)\bigr)
       -\frac{a^2}{2}(d-m)^2-3V.
\end{equation}
Inserting the variance expansion into \eqref{mid:parity-energy} and
using \eqref{mid:odd-cap} now gives
\begin{equation}\label{mid:energy-reduction}
 \mathcal E(d)\ge3v+\E\Psi-\frac{a^2}{2}\Var(d_o)
 \ge3v+\E\Psi-\frac r2 Q,
 \qquad Q=\frac{(1-r^2)w}{h^3}+\frac{r^2}{h}.
\end{equation}
Define
\begin{equation}\label{mid:Qstar}
 Q_* =\frac{2(2s-1)}{s(1+s)}
      =\frac6{1+s}-\frac2s.
\end{equation}
Since \(r=v(1+s)\), the identity
\(3v-rQ_*/2=r/s\) shows that the following strict inequality suffices:
\begin{equation}\label{mid:target}
 \E\Psi>\frac r2(Q-Q_*).
\end{equation}

\subsection{Pointwise bounds and the bias margins}

Near \(m=r^2/2\), a positive multiple of \(V^2\) supplies the needed
gain because \(\E V=v\). For \(m\ge1/5\), the gap
\(Q_*-Q\) will instead allow a small negative multiple of \(V\).
We first use these estimates to prove Proposition~\ref{prop:intermediate},
then prove the estimates in angular coordinates.

\begin{lemma}\label{mid:pointwise}
For every \(-1<d<1\), with \(V,\Psi\) defined by
\eqref{mid:V} and \eqref{mid:Psi}, one has
\begin{equation}\label{mid:pointwise-bounds}
 \begin{array}{ll}
 \Psi\ge(\frac12-2m^2)V^2,&0\le m\le\frac15,\\[2pt]
 \Psi\ge-\frac1{100}V,&\frac15\le m\le\frac23,\\[2pt]
 \Psi\ge-\frac1{50}V,&\frac23\le m\le\frac{37}{50}.
 \end{array}
\end{equation}
\end{lemma}

\begin{proof}[Proof of Proposition~\ref{prop:intermediate}]
Assume \(b=h-1+s\). Since \(s\ge21/25\),
\begin{equation}\label{mid:r-bound}
 r^2\le(184/625)^2<87/1000.
\end{equation}
The first term of the bound on \(w\), namely \(w\le1-m\), gives
\begin{equation}\label{mid:Qzero}
 Q\le Q_0:=\frac{1+r^2m}{h(1+m)}.
\end{equation}

First suppose \(r^2/2\le m\le1/5\). In particular \(m>0\).
On this interval,
\[
 h\ge1-\frac{51}{100}m^2,
\]
as follows by squaring: the difference of the squares is
\(m^2(1/50-2601m^2/10000)\ge0\). Consequently
\begin{align}
 \frac{1-Q_0}{m}
 &=\frac1h\left[\frac{1-r^2}{1+m}-\frac{1-h}{m}\right]
 \ge L(m):=\frac{913}{1000(1+m)}-\frac{51m}{100}.
 \label{mid:small-margin-left}
\end{align}
Here the bracket is positive, since it is at least
\(L(m)\ge L(1/5)=3953/6000>0\); thus the factor \(1/h\ge1\)
has been removed in the correct direction.

We compare \(L(m)\) with
\begin{equation}\label{mid:small-margin-right}
 R(m,s):=\frac{2v}{1+s}\left(\frac{v}{sm}+2m\right).
\end{equation}
For \(m\le3/50\), the hypothesis \(m\ge r^2/2\) gives
\[
 R(m,s)\le\frac4{s(1+s)^3}+\frac{4vm}{1+s}
 \le\frac4{(21/25)(46/25)^3}
     +\frac{4(4/25)(3/50)}{46/25}.
\]
For \(3/50\le m\le1/5\), use instead
\[
 R(m,s)\le
 \frac{2(4/25)^2}{(21/25)(46/25)(3/50)}
 +\frac{4(4/25)(1/5)}{46/25}.
\]
Since \(L\) decreases, the exact comparisons are
\[
\begin{array}{c|c|c}
 \text{range of }m & \text{lower bound for }L(m)
                    &\text{upper bound for }R(m,s)\\ \hline
 \rule{0pt}{22pt}(0,3/50] & \dfrac{220141}{265000}>\dfrac45
          & \dfrac{10032241}{12775350}<\dfrac45\\[7pt]
 [3/50,1/5]&\dfrac{3953}{6000}>\dfrac{16}{25}
           &\dfrac{4504}{7245}<\dfrac{16}{25}.
\end{array}
\]
Equations~\eqref{mid:small-margin-left}--\eqref{mid:small-margin-right}
therefore give \((1-Q)/m>R(m,s)\). Since
\[
 1-Q_* =\frac{v(2-s)}{s(1+s)},
\]
rearranging this strict inequality yields
\begin{equation}\label{mid:small-final-margin}
 \frac{1+s}{2v}(Q-Q_*)<\frac12-2m^2.
\end{equation}
The first line of Lemma~\ref{mid:pointwise}, Jensen's inequality, and
\eqref{mid:EV} give
\[
 \E\Psi\ge(\tfrac12-2m^2)\E V^2
           \ge(\tfrac12-2m^2)v^2
           >\frac r2(Q-Q_*),
\]
where the last step uses \(r=v(1+s)\) and
\eqref{mid:small-final-margin}.

Next suppose \(1/5\le m\le2/3\). For fixed
\(a_0=r^2\le87/1000\),
\[
 (\log Q_0)'=\frac{a_0}{1+a_0m}+\frac{2m-1}{1-m^2},\qquad
 (\log Q_0)''=-\frac{a_0^2}{(1+a_0m)^2}
       +\frac{2(m^2-m+1)}{(1-m^2)^2}>0.
\]
For the last sign, the positive term is at least \(3/2\), and
\(a_0^2<1\). Convexity bounds \(Q_0\) by its endpoint maximum;
it also increases with \(a_0\). At \(a_0=87/1000\), the two endpoint
values are
\[
 Q_0(1/5)=\frac{5087}{2400\sqrt6}<\frac{433}{500}=.866,
 \qquad
 Q_0(2/3)=\frac{4761}{2500\sqrt5}<\frac{213}{250}=.852.
\]
The strict comparisons follow respectively from
\[
 25435^2<6\cdot10392^2,\qquad 4761^2<5\cdot2130^2.
\]
Thus \(Q<.866\) throughout this range.

Finally suppose \(2/3\le m\le37/50\). The second term of the
bound on \(w\) gives
\begin{equation}\label{mid:Qhigh}
 Q\le R_0(m,a_0):=
 (1-a_0)\frac{\sqrt3}{(1+m)^{3/2}}+\frac{a_0}{\sqrt{1-m^2}},
 \qquad a_0=r^2\le87/1000.
\end{equation}
The coefficient of \(a_0\) is positive, because
\((1+m)^2>3(1-m)\) on this interval. Hence it suffices to take
\(a_0=87/1000\). Moreover,
\[
 \frac{\partial R_0}{\partial m}
 =-\frac{3(1-a_0)\sqrt3}{2(1+m)^{5/2}}
   +\frac{a_0m}{(1-m^2)^{3/2}}<0.
\]
Here is a rational comparison proving the last sign on the entire
interval. Since \(h>2/3\), the positive term is less than
\((87/1000)(37/50)(27/8)<1/4\). Since
\(\sqrt3>5/3\), \((1+m)^{5/2}<6\), and \(1-a_0\ge913/1000\),
the magnitude of the negative term is greater than
\(4565/12000>1/3\). At the left endpoint,
\[
 R_0(2/3,87/1000)=\frac{4761}{2500\sqrt5}<.852,
\]
so \(Q<.852\) in the final range.

To finish both remaining ranges, observe that
\[
 Q_*'(s)=\frac{2+4s(1-s)}{s^2(1+s)^2}>0,
 \qquad Q_*(21/25)=\frac{425}{483}>.879,
 \qquad \frac{1+s}{2}\ge\frac{23}{25}=.92.
\]
The exact positive margins after the respective pointwise losses are
\begin{align*}
 \frac{1+s}{2}(Q_*-Q)-\frac1{100}
 &>.92(.879-.866)-.01=\frac{49}{25000}>0,
       &&1/5\le m\le2/3,\\
 \frac{1+s}{2}(Q_*-Q)-\frac1{50}
 &>.92(.879-.852)-.02=\frac{121}{25000}>0,
       &&2/3\le m\le37/50.
\end{align*}
By the appropriate line of Lemma~\ref{mid:pointwise},
\(\E\Psi\ge-\eps\E V=-\eps v\); the last inequalities and
\(r=v(1+s)\) therefore imply \eqref{mid:target}.
All three ranges give \(\mathcal E(d)>r/s\) by
\eqref{mid:energy-reduction}, completing the proof.
\end{proof}

\subsection{Proof of the pointwise comparison}

\begin{proof}[Proof of Lemma~\ref{mid:pointwise}]
Introduce the midpoint and signed half-difference of the two angles:
\[
 \beta=\arcsin m,\qquad t=\frac{\arcsin d-\beta}{2},\qquad
 c=t+\beta.
\]
Then \(|c|+|t|=\max\{|\arcsin d|,|\beta|\}<\pi/2\).
The case \(t=0\) gives \(d=m\) and \(V=\Psi=0\); hence assume
\(t\ne0\). Elementary angle identities give
\[
 d-m=2\cos c\sin t,\qquad md+hH(d)=\cos(2t).
\]
Changing variables \(z=\sin\theta\) in the integral defining \(F\)
gives, with the orientation of the integral retained if \(t<0\),
\[
 g-F(m)=\int_{c-t}^{c+t}\frac{d\theta}{\sqrt{\cos\theta}}.
\]
Normalize this integral by setting
\[
 T=\frac mh,\qquad q=\frac{\cos c}{h}=\cos t-T\sin t>0,
 \qquad G=\tan c.
\]
It follows that
\begin{equation}\label{mid:angular-identities}
 V=\frac{2\sin^2t}{h},\qquad
 \frac{\Psi}{V}=4\sqrt q\,A-q^2-3,
 \qquad
 A=\frac{t}{\sin t}\frac12\int_{-1}^1
       (\cos(tu)-G\sin(tu))^{-1/2}\,du.
\end{equation}

To bound \(A\) from below, we keep the first three terms of a positive
expansion. If
\(y=\sin^2(tu)\), symmetry in \(u\) expresses the averaged integrand as
\[
 \sum_{j\ge0}\frac{\binom{4j}{2j}}{16^j}
     G^{2j}y^j(1-y)^{-j-1/4}.
\]
This expansion is convergent because \(|G\tan(tu)|<1\), and its
power coefficients in \(y\) are nonnegative. Keeping the terms through
degree two therefore gives the lower bound
\[
 1+\left(\frac14+\frac{3G^2}{8}\right)y
  +\left(\frac5{32}+\frac{15G^2}{32}
                 +\frac{35G^4}{128}\right)y^2.
\]
Concavity of sine on \([0,\pi/2]\) implies
\(\sin^2(tu)\ge u^2\sin^2t\). Averaging and writing
\(Y=\sin^2t\), we obtain
\[
 \frac12\int_{-1}^1(\cos(tu)-G\sin(tu))^{-1/2}\,du
 \ge1+\left(\frac1{12}+\frac{G^2}{8}\right)Y
       +\left(\frac1{32}+\frac{3G^2}{32}
                       +\frac{7G^4}{128}\right)Y^2.
\]
Multiply by the positive-series bound
\(t/\sin t\ge1+Y/6+3Y^2/40\), and discard terms of degree at
least three. The result is
\begin{equation}\label{mid:Taylor}
 \frac{4(A-1)}{\sin^2t}\ge P_0,
 \qquad
 P_0=1+\frac{G^2}{2}
 +\sin^2t\left(\frac{173}{360}+\frac{11G^2}{24}
                                  +\frac{7G^4}{32}\right).
\end{equation}
To express the remaining terms rationally, put \(x=\tan(t/2)\).
Then
\[
 q-1=-\sin t(T+x),\qquad
 4\sqrt q-q^2-3
 =-(q-1)^2\left(1+\frac2{(1+\sqrt q)^2}\right).
\]
Equations~\eqref{mid:angular-identities} and \eqref{mid:Taylor}
therefore imply
\begin{equation}\label{mid:bracket}
 \frac{\Psi}{V}\ge\sin^2t
 \left[\sqrt q\,P_0
 -(T+x)^2\left(1+\frac2{(1+\sqrt q)^2}\right)\right].
\end{equation}

We now bound these square-root coefficients on the required angular domains.

For \(m\le1/5\), \(T\le1/\sqrt{24}<21/100\) and
\(\beta<21/100\). For \(m\le37/50\), \(\beta<21/25\).
These bounds on \(\beta\) follow from \(\sin u\ge u-u^3/6\)
at the respective rational values. Moreover,
\[
 -\frac{\pi/2+\beta}{4}<\frac t2
 <\frac{\pi}{8},\qquad \tan(\pi/8)=\sqrt2-1<\frac{21}{50}.
\]
Using \(\pi<22/7\), the lower angle in absolute value is less than
\(9/20\) in the first range, and less than \(603/1000\) in the
second. The bounds
\[
 \tan(9/20)<1/2,\qquad \tan(603/1000)<7/10
\]
follow by dividing
\(\sin u\le u-u^3/6+u^5/120\) by
\(\cos u\ge1-u^2/2>0\); the resulting comparisons are rational.
Thus \(m\le1/5\) gives \(0\le T\le21/100\) and \(|x|\le1/2\),
whereas the two larger ranges have \(-7/10\le x\le21/50\).
The corresponding bounds on \(T\) are
\[
 T\le2/\sqrt5<9/10\quad(m\le2/3),\qquad
 T\le37/\sqrt{1131}<111/100\quad(m\le37/50).
\]

On the small rectangle,
\(q=(1-x^2-2Tx)/(1+x^2)\ge(27/50)/(5/4)=54/125=.432\),
so \(\sqrt q>13/20\). The exact identity
\begin{equation}\label{mid:root-identity}
 \sqrt q=1+\frac{q-1}{2}
          -\frac{(q-1)^2}{2(1+\sqrt q)^2}
\end{equation}
and the rational comparisons
\[
 \frac{200}{1089}<\frac{19}{100},\qquad
 \frac{1889}{1089}<\frac74
\]
give
\[
 \sqrt q\ge1+\frac{q-1}{2}-\frac{19}{100}(q-1)^2,
 \qquad 1+\frac2{(1+\sqrt q)^2}<\frac74.
\]

For the two larger ranges we use a bound on the actual angular domain:
\[
 q\ge\frac{\cos((\pi/2+\beta)/2)}h
   =\frac1{\sqrt{2(1+m)}}
   \ge\frac5{\sqrt{87}}>\frac{53}{100}.
\]
In particular \(\sqrt q>18/25\), and
\[
 \frac{625}{3698}<\frac{17}{100},\qquad
 \frac{3099}{1849}<\frac{42}{25}.
\]
Equation~\eqref{mid:root-identity} therefore gives
\[
 \sqrt q\ge1+\frac{q-1}{2}-\frac{17}{100}(q-1)^2,
 \qquad 1+\frac2{(1+\sqrt q)^2}<\frac{42}{25}.
\]

These quadratic lower bounds for \(\sqrt q\) make the bracket in
\eqref{mid:bracket} rational in \(T,x\). We clear its positive
denominators using the polynomials
\begin{align}
 &D=1+x^2,\qquad N=1-x^2-2Tx,\qquad
 z=2x,\qquad k=2x+T(1-x^2),
 \label{mid:polynomial-variables}\\
 &P=D^2\left(N^4+\frac{k^2N^2}{2}\right)
 +z^2\left(\frac{173N^4}{360}+\frac{11k^2N^2}{24}
                              +\frac{7k^4}{32}\right),
 \label{mid:P}\\
 &R_e=D^2+\frac{(N-D)D}{2}-e(N-D)^2.
 \label{mid:R}
\end{align}
At the angular parameters,
\begin{equation}\label{mid:positive-denominators}
 \begin{gathered}
 q=\frac ND,\qquad G=\frac kN,\qquad
 \sin t=\frac zD,\qquad P_0=\frac{P}{D^2N^4},\\
 D>0,\qquad N>0,\qquad P>0.
 \end{gathered}
\end{equation}
The needed polynomial signs are
\begin{align}
 \mathcal S(T,x)
 &:=R_{19/100}P-D^4N^4
       \left[1-\frac{33}{10}T^2+\frac74(T+x)^2\right]\ge0
 \label{mid:small-sign}\\[-3pt]
 &\hspace{35mm}(0\le T\le21/100,\ |x|\le1/2),\notag\\
 \mathcal L_{\eps}(T,x)
 &:=z^2\left[R_{17/100}P-\frac{42}{25}D^4N^4(T+x)^2\right]
       +\eps D^6N^4>0
 \label{mid:large-sign}\\[-3pt]
 &\hspace{10mm}(0\le T\le T_0,\ -7/10\le x\le21/50),\notag
\end{align}
where \((T_0,\eps)=(9/10,1/100)\) or \((111/100,1/50)\).
Their exact finite verification is given below.

For \(m\le1/5\), the lower bound for \(\sqrt q\) is
\(R_{19/100}/D^2\). Hence
\eqref{mid:small-sign} makes the bracket in \eqref{mid:bracket}
at least \(1-(33/10)T^2\).
To compare this with the claimed bound, put \(y=T^2\le441/10000\).
Both sides in the following squared comparison are nonnegative, and
\[
 \left(1-\frac{33}{10}y\right)^2(1+y)-(1-3y)^2
 =y\left(\frac25-\frac{471}{100}y+\frac{1089}{100}y^2\right)\ge0.
\]
Indeed \(2/5-(471/100)(441/10000)>0\). Consequently
\[
 1-\frac{33}{10}T^2
 \ge\frac{1-3T^2}{\sqrt{1+T^2}}
 =\frac{1-4m^2}{h}.
\]
Together with \(V=2\sin^2t/h\), this proves the first line of
\eqref{mid:pointwise-bounds}.

For the two larger ranges the lower bound for \(\sqrt q\) is
\(R_{17/100}/D^2\), and the coefficient of \((T+x)^2\) is at most
\(42/25\).
Since \(\sin^2t=z^2/D^2\), Equation~\eqref{mid:large-sign}
implies \(\Psi/V>-\eps\). This proves the other two lines for
\(t\ne0\), and their nonstrict form also holds at \(t=0\).
Notice that division by \(N\) has only been used at actual angular
points, where \eqref{mid:positive-denominators} holds. A larger
certification rectangle may contain \(N\le0\); its polynomial sign
assertion remains valid and is more than is needed here.
\end{proof}

\subsection{The six exact polynomial certificates}

For completeness, we specify the entire finite calculation behind
\eqref{mid:small-sign}--\eqref{mid:large-sign}. It uses rational
arithmetic and the tensor Bernstein basis described in
Appendix~\ref{app:arithmetic}. Explicitly, for
\(U(A,B)=\sum_{a,b}u_{ab}A^aB^b\) of coordinate degrees at most
\((n_1,n_2)\), its coefficients in that basis are
\begin{equation}\label{mid:Bernstein-transform}
 \mathcal B_{ij}(U;n_1,n_2)
 =\sum_{a=0}^{i}\sum_{b=0}^{j}
 u_{ab}\frac{\binom ia\binom jb}{\binom{n_1}a\binom{n_2}b}
 \quad(0\le i\le n_1,\ 0\le j\le n_2).
\end{equation}
The Bernstein basis functions are nonnegative on the unit square and
sum to one. Thus a positive minimum of these coefficients proves
positivity throughout that square.

For \(\mathcal S\), use \((A,B)\in[0,1]^2\), \(\sigma\in\{-1,1\}\),
and the two charts
\begin{align}
 C_{1,\sigma}(A,B)&=(21\sqrt A/100,\ \sigma B\sqrt A/2),
 \label{mid:chart-one}\\
 C_{2,\sigma}(A,B)&=(21B\sqrt A/100,\ \sigma\sqrt A/2).
 \label{mid:chart-two}
\end{align}
They cover the required rectangle: choose the first chart when
\(T/(21/100)\ge |x|/(1/2)\), and the second when the reverse
inequality holds. In each chart \(\mathcal S\circ C_{j,\sigma}\)
has a factor \(A\). To make the expansion entirely explicit, if
\(\mathcal S(T,x)=\sum c_{ij}T^ix^j\), then every nonzero term has
\(i+j\) even and at least two, and the divided polynomials are
\begin{align}
 U_{1,\sigma}
 &=\sum_{i,j}c_{ij}(21/100)^i(\sigma/2)^j
      A^{(i+j)/2-1}B^j,\label{mid:chart-expansion-one}\\
 U_{2,\sigma}
 &=\sum_{i,j}c_{ij}(21/100)^i(\sigma/2)^j
      A^{(i+j)/2-1}B^i.\label{mid:chart-expansion-two}
\end{align}
Their degrees are at most \((8,18)\) and \((8,6)\), respectively.
The values at \(A=0\) are polynomial extensions; after multiplying
back by \(A\ge0\), the origin is included in the nonstrict sign claim.

For \(\mathcal L_\eps\), subdivide into sixteen equal rectangles.
For \(j,k\in\{0,1,2,3\}\), use the affine map
\begin{equation}\label{mid:large-map}
 T=\frac{T_0}{4}(j+A),\qquad
 x=-\frac7{10}+\frac7{25}(k+B),\qquad (A,B)\in[0,1]^2.
\end{equation}
Each resulting polynomial has coordinate degrees at most \((6,20)\).
If the untransformed coefficient of \(T^ix^j\) is \(c_{ij}\), the
coefficient of \(A^aB^b\) after any affine map
\((T,x)=(\ell_T+w_TA,\ell_x+w_xB)\) is exactly
\begin{equation}\label{mid:affine-coefficients}
 \sum_{i\ge a,\ j\ge b}c_{ij}\binom ia\binom jb
 \ell_T^{i-a}w_T^a\ell_x^{j-b}w_x^b.
\end{equation}
Thus \eqref{mid:polynomial-variables}--\eqref{mid:R},
\eqref{mid:small-sign}--\eqref{mid:large-sign}, and
\eqref{mid:Bernstein-transform}--\eqref{mid:affine-coefficients}
specify every coefficient by a finite rational sum.

Table~\ref{mid:certificate-table} records the minimum coefficients.
An entry \(M\) in the last column means
\(\lfloor C\min\mathcal B_{ij}\rfloor=M\), with the minimum also
taken over all sixteen boxes in the last two rows. In particular the
minimum is at least \(M/C>0\). These six integer comparisons are
the finite arithmetic part of the pointwise proof; the accompanying
exact reconstruction scripts implement the displayed sums.

\begin{table}[htbp]
\centering
\begin{tabular}{llllr}
\toprule
Polynomial & Domain map & Orders & Scale \(C\) & \(M\)\\
\midrule
\(U_{1,-}\) & \eqref{mid:chart-one}, \(\sigma=-1\) & \((8,18)\) & \(1000\) & 90\\
\(U_{1,+}\) & \eqref{mid:chart-one}, \(\sigma=1\) & \((8,18)\) & \(1000\) & 27\\
\(U_{2,-}\) & \eqref{mid:chart-two}, \(\sigma=-1\) & \((8,6)\) & \(1000\) & 293\\
\(U_{2,+}\) & \eqref{mid:chart-two}, \(\sigma=1\) & \((8,6)\) & \(1000\) & 120\\
\(\mathcal L_{1/100}\) & \eqref{mid:large-map}, \(T_0=9/10\) & \((6,20)\) & \(10000\) & 53\\
\(\mathcal L_{1/50}\) & \eqref{mid:large-map}, \(T_0=111/100\) & \((6,20)\) & \(10000\) & 26\\
\bottomrule
\end{tabular}
\caption{Exact Bernstein certificate minima for the intermediate band.}
\label{mid:certificate-table}
\end{table}

